% Figure: the temperature-vs-heat-added profile of water carried through a
% once-through boiler at constant pressure. The fluid warms as subcooled
% liquid (sloped), then boils across the two-phase region at constant
% saturation temperature (flat plateau, the latent-heat leg), then superheats
% as dry vapour (sloped again). The flat plateau is the largest single block
% of heat and pins the pinch point in a heat exchanger.
\begin{figure}[htbp]
\centering
\begin{tikzpicture}
\begin{axis}[
  width=12cm, height=7.5cm,
  xlabel={heat added per kilogram $q$ (\si{\kilo\joule\per\kilogram})},
  ylabel={temperature $T$ (\si{\degreeCelsius})},
  xmin=0, xmax=3200, ymin=0, ymax=560,
  axis lines=left,
  xtick={0,632,2417,3199},
  xticklabels={$0$,$q_{\mathrm{sat}}$,$q_{\mathrm{sat}}+h_{fg}$,$q_{\mathrm{in}}$},
  ytick={0,100,200,295,400,500},
  yticklabels={$0$,$100$,$200$,$T_{\mathrm{sat}}$,$400$,$500$},
  tick label style={font=\scriptsize},
  label style={font=\small},
  clip=true,
]
% subcooled-liquid warming: 46 C at q=0 up to 295 C (sat) at q=632
\addplot[engblue, very thick] coordinates {(0,46) (632,295)};
% two-phase boiling plateau at 295 C, q from 632 to 632+1785=2417
\addplot[engred, very thick] coordinates {(632,295) (2417,295)};
% superheating dry vapour: 295 C at q=2417 up to 500 C at q=3199
\addplot[engorange, very thick] coordinates {(2417,295) (3199,500)};
% saturation-temperature guide line
\addplot[enggray, dashed, thick] coordinates {(0,295) (3199,295)};
\node[anchor=south west, font=\scriptsize, engblue] at (axis cs:20,120)
  {liquid warming};
\node[anchor=south, font=\scriptsize, engred] at (axis cs:1500,300)
  {boiling: constant $T_{\mathrm{sat}}$};
\node[anchor=west, font=\scriptsize, engorange] at (axis cs:2500,430)
  {superheat};
% endpoint markers
\addplot[engblack, only marks, mark=*, mark size=1.5pt] coordinates {
  (0,46) (632,295) (2417,295) (3199,500)
};
\end{axis}
\end{tikzpicture}
\caption[Temperature against heat added through a boiler at constant
pressure]{Water carried through an $8\,\si{\mega\pascal}$ boiler at constant
pressure warms in three legs. As subcooled liquid it warms steeply, because the
liquid's heat capacity turns each kilojoule into a large temperature rise; at
the saturation temperature it boils across a flat plateau, absorbing the latent
heat $h_{fg}$ at constant temperature; then as dry vapour it superheats along a
third sloped leg. The flat plateau is the largest single block of heat and the
reason a boiler spends most of its length at one temperature. The horizontal
plateau also fixes the pinch point of any counter-flow heat exchanger boiling a
fluid: the hot stream must stay above the flat saturation temperature across
the whole boiling leg, which is the geometric constraint the Rankine
state-point analysis and the two-phase heat exchangers of the volume design
around.}
\label{fig:vol04ch02:boiler-tq-profile}
\end{figure}
